% !TEX TS-program = lualatex

% ------------------------------------------------- %
% -- DO NOT MODIFY. FILE GENERATED AUTOMATICALLY -- %
% ------------------------------------------------- %

\documentclass{tutodoc}

\usepackage{../preamble.cfg}


\begin{document}

% -- AUTOMATICALLY GENERATED UGLY CODE - START -- %

\subsection{Lua}
\label{products-lua}

\begin{tdocnote}
Initially, the {\thisproj} project was created to provide ready-to-use color palettes for {\LUADRAW}, a package that greatly facilitates the creation of high-quality 2D and 3D plots using {\LuaLaTeX} and {\TIKZ}. The {\lua} implementation is now integrated into {\LUADRAW}.
\end{tdocnote}

\subsubsection{Basic use}

The {\lua} palette names all use the prefix \tdoccodein{text}+pal+ followed by the name available in the file \tdoccodein{text}+palettes.json+. You can access a palette by two ways.

\begin{itemize}
\item
\tdoccodein{lua}+palGistHeat+ is a {\lua} variable.

\item
\tdoccodein{lua}+getPal('GistHeat')+ and \tdoccodein{lua}+getPal('palGistHeat')+ are equal to \tdoccodein{lua}+palGistHeat+.
\end{itemize}

The {\lua} palette variables are arrays of ten arrays of three floats (making it straightforward to use a color from a palette). For example, the definition of \tdoccodein{lua}+palGistHeat+ looks like the following partial code.

\begin{tdoccode}{lua}
palGistHeat = {
    {0.0, 0.0, 0.0},
    {0.105882, 0.0, 0.0},
    {0.211764, 0.0, 0.0},
    -- ... With 7 more RBG colors.
}
\end{tdoccode}

\subsubsection{Creating palettes from existing ones}

The \tdoccodein{lua}+getPal+ function provides options to build new palettes by transforming existing ones. The following example shows how to do this (all options are used).

\begin{tdoccode}{lua}
BlackbodyTransformed = getPal(
    'Blackbody',
    {
        extract = {2, 5, 8, 9},
        shift   = 1,
        reverse = true
    }
)
\end{tdoccode}

% -- AUTOMATICALLY GENERATED UGLY CODE - END -- %

\end{document}
